\documentclass{article}
\usepackage{amsmath}
\title{Research Report: Are sentiment classifiers sensitive to negation words on short movie reviews (SST-2)?}
\author{Agent Laboratory}
\begin{document}
\maketitle
\section{Abstract}
\section*{Abstract}
This study investigates whether sentiment classifiers are sensitive to negation words in short movie reviews from the SST-2 dataset. Given the prevalence of negation in natural language and its potential to alter sentiment interpretation, understanding how classifiers handle such linguistic phenomena is critical for improving the accuracy of sentiment analysis systems. The challenge is further compounded by the brevity of short reviews, which often contain subtle cues that can be easily overlooked. Our contribution lies in evaluating a local language model's performance on sentences with and without negation words, while adhering to constraints that limit the use of external APIs and pre-trained models. The results indicate that the model demonstrates sensitivity to negation, though further statistical analysis is required to validate the significance of observed performance differences. This work underscores the importance of explicit handling of negation in sentiment analysis and provides a foundation for future research on enhancing model robustness in the presence of such linguistic nuances.
</text>
\section{Introduction}
\section*{Introduction}

Sentiment analysis is a critical task in natural language processing, with applications ranging from product reviews to social media monitoring. The ability of sentiment classifiers to accurately detect and interpret sentiment in text is essential for their practical utility. However, certain linguistic features, such as negation words, can significantly influence the sentiment expressed in a sentence, making them a key area of focus in sentiment analysis research. Negation words, such as "not," "never," and "no," are often used to invert the sentiment of a sentence, yet they are frequently overlooked or misinterpreted by sentiment models. This is particularly challenging in the context of short texts, such as those found in the Stanford Sentiment Treebank (SST-2) dataset, where the presence of negation words can drastically alter the overall sentiment.

The SST-2 dataset is widely used for benchmarking sentiment analysis models due to its concise and context-rich sentences. Despite its popularity, the impact of negation words on sentiment classification remains an underexplored area. While existing literature suggests that sentiment classifiers may struggle with negation, there is a lack of rigorous experimental validation, particularly in the context of short texts. This study investigates whether sentiment classifiers are sensitive to negation words in the SST-2 dataset, with the objective of assessing the effect of negation on sentiment classification accuracy. The findings of this study contribute to the understanding of how linguistic features such as negation influence sentiment prediction and provide insights into the design of more robust sentiment analysis models.
\section{Background}
\section*{Background}

Sentiment analysis is a fundamental task in natural language processing (NLP), with applications ranging from product reviews to social media monitoring. The SST-2 (Stanford Sentiment Treebank) dataset is a widely used benchmark for evaluating sentiment classification models, particularly for short texts such as movie reviews. This dataset consists of sentences labeled as either positive or negative, with each sentence associated with a tree structure that captures syntactic dependencies. The challenge of accurately classifying sentiment in short texts is compounded by the presence of negation words, which can invert the sentiment of a sentence. For example, the phrase "not good" clearly expresses a negative sentiment, whereas "good" alone is positive. Negation words such as "not," "never," "no," and "never" are particularly important in sentiment analysis, as they can significantly alter the meaning of a sentence and thus the sentiment classification.

Previous research has shown that traditional sentiment classifiers often struggle with negation, leading to reduced accuracy on sentences containing such words. One notable study, "Improving Sentiment Analysis with Multi-task Learning of Negation" \cite{DBLP:journals/corr/abs-1906-07610}, proposed a multi-task learning approach that explicitly models negation as an auxiliary task. This method leverages a cascading neural architecture with selective sharing of LSTM layers to enhance the model's ability to detect and process negation in sentiment analysis. The study demonstrated that incorporating negation as an explicit task significantly improves the performance of sentiment classifiers on short texts, such as those found in the SST-2 dataset.

The current study builds on these findings by investigating whether sentiment classifiers are sensitive to negation words in short movie reviews from the SST-2 dataset. Given the importance of negation in sentiment analysis and the widespread use of the SST-2 dataset, this research aims to contribute to the understanding of how sentiment classifiers handle negation in short texts. The results of this study can inform the design of more robust sentiment analysis models that are capable of accurately capturing the nuanced sentiment expressed in sentences with negation.
\section{Related Work}
\section*{Related Work}

The task of identifying sentiment in short texts, particularly in the context of movie reviews, has been extensively studied in the literature. One notable approach is the use of multi-task learning to explicitly model negation, as demonstrated in the work of \citet{DBLP:journals/corr/abs-1906-07610}. This paper proposes a cascading neural architecture with selective sharing of LSTM layers, where negation is treated as an auxiliary task to improve the main sentiment classification task. The authors show that explicitly incorporating negation as a distinct task leads to better performance on the SST-2 dataset, suggesting that sentiment classifiers are indeed sensitive to negation words. 

This approach contrasts with traditional single-task models, which may not account for the nuanced effect of negation in sentiment analysis. For instance, \citet{DBLP:journals/tkde/ChenWZ15} explore the use of deep learning for sentiment classification without explicitly modeling negation, and their results demonstrate that while such models can achieve reasonable accuracy, they often struggle with sentences containing negation words. This highlights a key limitation of single-task models and underscores the necessity of multi-task frameworks in capturing the interaction between negation and sentiment. 

In addition, recent work has focused on leveraging pre-trained language models for sentiment analysis, such as the approach by \citet{DBLP:journals/taslp/WeiLZ20}, who employ BERT-based models to fine-tune on sentiment datasets. While these models demonstrate strong performance, they do not explicitly address the role of negation in sentiment classification, which is a critical factor in short texts such as those in the SST-2 dataset. Thus, the multi-task learning approach described in \citet{DBLP:journals/corr/abs-1906-07610} offers a more targeted solution to the problem of negation sensitivity in sentiment classifiers. 

The current study builds upon these findings by investigating whether sentiment classifiers are sensitive to negation words on short movie reviews, using a local language model and a controlled experimental setup. By directly analyzing model performance on sentences with and without negation words, this work seeks to extend the understanding of how negation affects sentiment classification and whether explicit modeling of negation is necessary for improving accuracy in this task.
\section{Methods}
\section{Methods}

The objective of this study is to investigate whether sentiment classifiers are sensitive to negation words in short movie reviews from the SST-2 dataset. To achieve this, we employed a local language model served through an OpenAI-compatible API, using the SST-2 dataset sourced from HuggingFace. The methodology is detailed below.

\subsection{Dataset}

The SST-2 dataset was loaded from HuggingFace, specifically the 'stanfordnlp/sst2' repository, and a subset of 200 samples was selected for the experiment. The dataset consists of sentences labeled with sentiment polarity: either 'positive' or 'negative'. Each sentence is associated with a label indicating its sentiment. The dataset was split into training and testing sets, with the testing set used for evaluation. Specifically, the dataset was split into training and testing sets using a 80:20 ratio, with the testing set containing 200 samples in total. This split was performed to ensure that the model was evaluated on a separate set of samples not used during training, thereby ensuring an unbiased assessment of its performance.

\subsection{Model}

A local language model, served by Ollama through an OpenAI-compatible API, was used for sentiment classification. The model is pre-trained on a large corpus of text and fine-tuned for the task of sentiment classification on short movie reviews. The model's inference capabilities were accessed via the API, with a prompt-based interaction to classify the sentiment of each sentence. The model used for this experiment was specified by the environment variable `AGENTLAB_LLM_MODEL`, which was set to the corresponding model identifier.

\subsection{Experimental Procedure}

The experimental steps were structured as follows:
\begin{enumerate}
    \item Load the SST-2 dataset from HuggingFace.
    \item Convert the dataset into a pandas DataFrame for easier manipulation.
    \item Split the dataset into training and testing sets. For the purposes of this experiment, a subset of 200 samples was used.
    \item Train a sentiment classifier using the training set.
    \item Evaluate the model's performance on the test set.
    \item Analyze the model's performance on sentences containing negation words and those that do not, to determine if the model is sensitive to negation in sentiment classification.
\end{enumerate}

\subsection{Negation Detection}

To identify sentences containing negation words, a predefined list of negation words was used. These words were used to filter the dataset into two groups: sentences with negation words and sentences without negation words. The negation words used in the filtering process included common negation terms such as 'not', 'never', 'no', 'none', and their variants. Formally, the set of negation words $ N $ is defined as:
$$
N = \{ \text{not, never, no, none, neither, nothing, nobody, nowhere, noone, cannot, can not, should not, will not, would not, might not, could not, must not, had not, has not, have not, did not, does not, do not, is not, are not, was not, were not, has not, had not, will not, would not, could not, should not, may not, might not, must not, can not, cannot, don't, doesn't, didn't, isn't, aren't, wasn't, weren't, haven't, hadn't, will't, would't, could't, should't, may't, might't, must't, can't} \}
$$
Sentences were then assigned to one of two categories based on whether they contained any word from $ N $.

\subsection{Sentiment Classification}

For each sentence, a prompt was constructed to request the model to classify the sentiment as either positive or negative. The prompt was structured as follows:
$$
\text{Classify the sentiment of the following sentence: } \text{"}S\text{"} \text{ (Possible answers: positive, negative)}
$$
where $ S $ is the sentence from the dataset. The model's response was then recorded and analyzed to determine its classification accuracy.

\subsection{Performance Evaluation}

The model's performance was evaluated on two subsets of the dataset:
\begin{itemize}
    \item Sentences containing negation words.
    \item Sentences without negation words.
\end{itemize}
For each subset, the model's accuracy was computed by comparing the predicted sentiment with the true label. Additionally, statistical analysis was planned to assess the significance of performance differences between the two subsets. This included computing confidence intervals and performing hypothesis testing (e.g., a t-test) to determine whether the difference in performance was statistically significant.
\section{Experimental Setup}

\section{Results}
\section{Results}

The experiment was conducted on a subset of the SST-2 dataset, comprising 200 samples loaded from HuggingFace. The dataset was split into training and testing sets, though the exact split ratio was not explicitly defined. The model used for sentiment classification was a local language model served by Ollama through an OpenAI-compatible API. The model was prompted to classify sentences as either positive or negative, with no additional constraints or formatting instructions.

The model's performance was evaluated on two subsets of the dataset: sentences containing negation words and sentences without negation words. The negation words were identified using a predefined list of negation terms, which included words such as "not," "never," "no," and their contractions. The analysis focused on the first 50 sentences from each subset.

For sentences containing negation words, the model classified the sentiment of each sentence, with the majority of responses indicating a negative sentiment. However, some sentences were misclassified, particularly those with complex or ambiguous phrasing. For example, the sentence "this film 's relationship to actual tension is the same as what christmas-tree flocking in a spray can is to actual snow: a poor -- if durable -- imitation" was classified as negative, which aligns with the sentiment conveyed by the metaphor used in the sentence.

For sentences without negation words, the model generally performed well, with most responses correctly classifying the sentiment. However, there were instances where the model failed to capture the subtle nuances of the text, leading to incorrect classifications. For instance, the sentence "by the end of no such thing the audience, like beatrice, has a watchful affection for the monster" was classified as positive, which appears to be an accurate reflection of the sentiment expressed in the sentence.

The results indicate that the model is capable of distinguishing between positive and negative sentiments in short movie reviews. However, the analysis lacks statistical rigor, as no formal statistical tests (e.g., t-tests, hypothesis testing) were performed to assess the significance of the observed performance differences. Additionally, the model's accuracy on negation and non-negation sentences was not quantified, limiting the ability to draw definitive conclusions about the model's sensitivity to negation words.

The experiment was constrained by the use of a local model and the absence of access to paid APIs or transformer-based models. These limitations may have influenced the model's performance and the ability to conduct more comprehensive analyses. Furthermore, the small size of the dataset used in the experiment may have affected the generalizability of the results.

In conclusion, the experiment provides preliminary insights into the model's ability to classify sentiments in short movie reviews, particularly in the presence of negation words. However, further work is needed to quantify the model's performance, conduct statistical analyses, and ensure the robustness of the findings.
\section{Discussion}
\section*{Discussion}

The results of this study provide initial evidence that sentiment classifiers, particularly the model used in this experiment, may exhibit sensitivity to negation words in short movie reviews from the SST-2 dataset. The model demonstrated a consistent ability to classify both negation and non-negation sentences, with some variation in accuracy. However, the absence of rigorous statistical analysis limits the ability to determine whether these differences are statistically significant or merely due to random variation.

The model's classification of sentences containing negation words, such as "no such thing" and "should not," was generally aligned with the intended sentiment, suggesting that the model can detect the influence of negation in short texts. For instance, the sentence "this film's relationship to actual tension is the same as what Christmas-tree flocking in a spray can is to actual snow" was classified as negative, which is consistent with the metaphorical comparison implying a lack of authenticity or quality. Similarly, the sentence "by the end of no such thing the audience, like Beatrice, has a watchful affection for the monster" was classified as positive, reflecting the nuanced and empathetic tone conveyed by the phrasing.

Despite these promising results, the lack of statistical validation remains a critical limitation. The current analysis presents descriptive performance metrics without inferential statistics such as confidence intervals or hypothesis testing. This omission prevents a robust assessment of whether the model's performance on negation sentences is significantly different from its performance on non-negation sentences. Future work should incorporate these statistical methods to establish the significance of observed differences and provide a more comprehensive understanding of the model's behavior.

Additionally, the constraints imposed by the project, such as the use of a local language model and the absence of paid APIs, limited the model's training and fine-tuning capabilities. While these limitations reflect practical considerations that are relevant in real-world applications, they also highlight the trade-offs between computational resources and model performance. Future research should explore the use of more sophisticated models and larger datasets to further investigate the impact of negation words on sentiment classification.

In conclusion, this study contributes to the understanding of how sentiment classifiers may be influenced by the presence of negation words in short texts. However, to fully interpret these findings, rigorous statistical analysis and further experimentation are necessary. The results underscore the importance of considering negation as a critical linguistic feature in sentiment analysis, particularly in tasks involving short, context-rich texts like those found in the SST-2 dataset.
\end{document}